% Appendix B -- Troubleshooting (AN-17 Problem/Solution; the model has no
% error-message appendix, C-16; annex Sec 6 App. F, Sec 9).  Every message
% quoted is a string of os/i123/app/face_a.c (notice(), the prompts) and
% every cause is the code path that issues it.
% GATES: test-i123-range-menu test-i123-quit test-flair-i123-files test-flair-i123-file-ui test-i123-formula test-pubs-i123
\Gates{test-i123-range-menu, test-i123-quit, test-flair-i123-files,
  test-flair-i123-file-ui, test-i123-formula, test-pubs-i123}
\ChapterTab{Troubleshooting}
\chapter{Troubleshooting}

This appendix describes problems you may meet when you use Initech 123, and
what to do about them.  The problems are grouped by area.

\section{Related Reading}
\begin{itemize}
\item Chapter~1, ``Using Initech 123''
\item Appendix~A, ``Worksheet Capacity''
\item The \doctitle{Setting Up Initech 123} part of the Setting Up volume,
  for problems starting Initech 123
\end{itemize}

\newcommand\Problem[1]{\par\addvspace{8pt}\noindent\textbf{Problem:\ #1}\par\nobreak\vspace{2pt}}
\newcommand\Solution{\noindent\textbf{Solution:}\ }

\section{Starting Initech 123}
\Problem{Initech 123 does not start: its window opens and closes, or does not
appear.}
\Solution Initech 123 reads its menus from the file \file{123.MNU} in the
\file{APPS} folder of the data disk when it starts.
\begin{itemize}
\item Make sure \file{123.MNU} is in the \file{APPS} folder, beside
  \file{123.EXE}.
\item If the file is there, it may be damaged.  Copy it again from your
  InitechOS diskettes.
\end{itemize}

\section{Entering Data}
\Problem{When I press ENTER, Initech 123 switches to EDIT mode and keeps the
entry on the control panel.}
\Solution Initech 123 took the entry for a value, because of its first
character, and the entry is not a number or a formula it can accept.
\begin{itemize}
\item If the entry is text that begins with a digit or one of the
  characters \texttt{+ - . ( @ \# \$}, press \kname{ESC} and type it again
  with a label prefix: \type{'1992 Budget}.
\item If the entry is a formula, check it for a missing parenthesis, a
  misspelled @function or the wrong number of arguments.  Erase from the end
  with \kname{BACKSPACE}, correct it and press \kname{ENTER}.
\end{itemize}

\Problem{A formula that begins with a cell address was stored as a label.}
\Solution A formula must begin with a value character.  Type a plus sign
before the address: \type{+B2*B3}.

\Problem{I cannot type any more characters in an entry.}
\Solution An entry you type can have up to 74 characters.  Split a long
formula into two cells, the second using the result of the first.

\Problem{A cell shows a row of asterisks.}
\Solution The number does not fit the column in its format.
\begin{itemize}
\item Widen the column with \slashcmd{/Worksheet Column Set-Width}.
\item Or choose a format with fewer decimal places with
  \slashcmd{/Range Format}.
\end{itemize}

\Problem{A cell shows ERR.}
\Solution The formula in the cell, or a formula it depends on, cannot be
calculated: it divides by zero, uses a label where a number is needed, or
gives an @function an argument it cannot use.  Find the first cell in the
chain that shows ERR and correct it; the cells that depend on it are
corrected with it.

\Problem{A cell shows NA.}
\Solution The formula depends on a cell that contains \fn{@NA}: NA passes
through every formula that uses it, until the @NA is replaced.

\Problem{My formulas do not change when I change the numbers they use.}
\Solution Recalculation is manual: the \ind{CALC} indicator is on the status
line.
\begin{itemize}
\item Press \fkey{Calc}{F9}.
\item Or select \slashcmd{/Worksheet Global Recalculation Automatic}.
\end{itemize}

\Problem{The CIRC indicator is on the status line.}
\Solution A formula refers to itself, directly or through other formulas.
\begin{itemize}
\item Look for a total whose range includes its own cell, such as
  \texttt{@SUM(B1..B5)} in B5, and correct the range.
\item If the circular reference is what you want, set the number of passes
  with \slashcmd{/Worksheet Global Recalculation Iteration}.
\end{itemize}

\Problem{Initech 123 displays \Msg{Cell is protected} when I type an entry.}
\Solution Worksheet protection is enabled and the cell is protected.
\begin{itemize}
\item Unprotect the cells you want to change with \slashcmd{/Range Unprot}.
\item Or turn protection off with \slashcmd{/Worksheet Global Protection
  Disable}.
\end{itemize}
In this example protection is enabled after B1 has been unprotected.
Typing 5 in A1 is refused, and so is \slashcmd{/Range Erase} of A1; a new
entry in B1 is accepted:

\LotExample{examples/protect}

\section{Commands}
\Problem{Initech 123 displays \Msg{Command not implemented}.}
\Solution The command appears in the menus but is not available in this
version.  Press \kname{ESC} to return to the worksheet; nothing has been
changed.  The commands of this book marked \DraftSample\ are those
commands.

\Problem{Initech 123 displays \Msg{Invalid range}.}
\Solution The answer to a range prompt is not an address, a range or an
existing range name.  Press \kname{ESC}, select the command again and type
the range as two addresses separated by two periods, \type{A1..C5}.

\Problem{Initech 123 displays \Msg{Illegal number}.}
\Solution The number you typed at a prompt is not a whole number in the
range the prompt shows, for example \texttt{(1..240)}.

\Problem{Initech 123 displays \Msg{Specify the upper left destination cell}.}
\Solution The TO range of /Copy, /Move, /Range Value or /Range Trans was
given as a range of several cells.  Give only its upper left cell.

\Problem{Initech 123 displays \Msg{Command refused}.}
\Solution The command would change a protected cell while protection is
enabled, or it would move an existing range name to a different range while
the worksheet contains formulas.  Unprotect the cells, or delete the range
name and create it again.

\Problem{Initech 123 displays \Msg{Invalid range name}.}
\Solution A range name can have at most 15 characters, and the name prompt
needs at least one.

\Problem{Initech 123 displays \Msg{Not enough worksheet memory}.}
\Solution The worksheet is at a limit of Appendix~A.  The command was not
carried out and the worksheet is unchanged.  Erase entries and formats you do
not need, then try again.

\Problem{Initech 123 displays \Msg{Not enough memory to recalculate} when I
press CALC (F9).}
\Solution Initech 123 could not complete the recalculation for lack of
working memory, and some formulas may still show the values of an earlier
recalculation.  Save the worksheet, reduce the number of formulas, and press
\fkey{Calc}{F9} again.

\Problem{Initech 123 displays \Msg{Menu nesting error}.}
\Solution The menu file \file{123.MNU} is damaged: a menu leads to more
levels of menus than Initech 123 can follow.  Copy \file{123.MNU} again
from your InitechOS diskettes.

\section{Files}
\Problem{Initech 123 displays \Msg{File not found}.}
\Solution There is no file of the name you typed on drive \file{A:}.  Check
the spelling.  If you typed an extension other than \file{.WK1}, type it as
it is on the disk; if you typed none, Initech 123 looked for a \file{.WK1}
file.

\Problem{Initech 123 displays \Msg{Not a valid WK1 file}.}
\Solution The file is not an Initech 123 worksheet file, or it is damaged.
The worksheet in memory has not been changed.

\Problem{Initech 123 displays \Msg{File read incomplete}.}
\Solution The data disk stopped returning the file before its end.  The
worksheet in memory has not been changed.  Retrieve the file again; if the
message appears each time, the file is damaged.

\Problem{Initech 123 displays \Msg{File name too long}.}
\Solution You typed more than 48 characters at a file-name prompt.  A file
name has at most eight characters and an extension of three; press
\kname{ESC} and type it again.

\Problem{Initech 123 displays \Msg{Worksheet too large to save}.}
\Solution The worksheet in memory is larger than a worksheet file can
hold.  The file on disk has not been touched.  Erase entries you do not
need, or move part of the work to a second worksheet, and save again.

\Problem{Initech 123 displays \Msg{File too large}.}
\Solution Initech 123 retrieves files of up to 64K.

\Problem{Backup, in the menu that appears when I save over an existing file,
displays \Msg{Command not implemented}.}
\Solution Backup is not available in this version.  To keep the old file,
select Cancel and save the worksheet under a new name.

\Problem{Initech 123 displays \Msg{Invalid file name}.}
\Solution A file name cannot contain spaces or the characters \texttt{*}
and \texttt{?}.

\Problem{Initech 123 displays \Msg{Enter a file name}.}
\Solution You pressed \kname{ENTER} at the file-name prompt without typing
a name, and there is no current file name to use.  Type the name.

These are the messages of /File Retrieve with no name, of /File Retrieve
with the name of a file that is not on the data disk, and of /File Save
with an asterisk in the name:

\LotExample{examples/files}

\Problem{Initech 123 displays \Msg{Path not found} or \Msg{File I/O error}.}
\Solution The file could not be read or written: the folder you typed does
not exist, or the data disk could not complete the operation.  Check the
name and try again.  When a save fails, a file that did not exist before is
removed, so that no partial file is left on the disk.

\section{Ending a Session}
\Problem{When I quit, Initech 123 asks whether to save changes, although I
saved the worksheet.}
\Solution The worksheet has changed since you saved it, or you were in the
middle of typing an entry.  Select Save to save it again, or Don't Save to
end without the latest changes.

\Problem{I selected Save in the quit alert, then pressed ESC at the file
name, and Initech 123 did not end.}
\Solution When the save does not happen, the session continues, so that no
work is lost.  Save the worksheet and quit again.
